\section{Task and Evaluation}
\label{sec:task}
\paragraph{Semantic target.}
Q1 evaluates six objects in simulation: a marked cube, a wooden die, a duck, a hedgehog, a pig, and a milk carton. Q2 and Q3 use only the marked cube. Instructions request one of six surface letters, or one through six dots on the die. The cube letters are $\{T,E,R,L,S,H\}$. Each target has a fixed reference orientation $Q(g)$ that places the requested face upward. Scoring compares the full quaternion, including rotation about the face normal. This rule is shared by all methods.

At step $t$, the agent receives an RGB image $I_t$ and instruction $\ell$. The motor policy receives proprioception $x_t$ and measured pose $\hat p_t$. Cube pose is estimated from RGB; other objects use noisy measurements constructed from simulator pose. Controllers do not directly receive the unperturbed true pose $p_t$. Teacher actions and privileged targets are used for training, and true pose is used for independent scoring. During deployment, the controller cannot use a future agent reply. We denote the active goal vector by $g_t$, the recurrent state by $h_t$, and a pending request by $m_t=1$.

\paragraph{Timing.}
Only one agent request can be pending at a time. Cube episodes allow up to six model calls and 3,600 control steps at 30~Hz. Other objects use 20~Hz over the same 120 seconds of simulated time. For the cube, a call taking $L$ wall-clock seconds can affect control only after it completes and at least $\lceil30L\rceil$ steps have passed since the request. Physics continues while the agent reasons. The simulator runs slower than real time, so this models inference delay but does not demonstrate a 30~Hz hardware system.

\paragraph{Autonomous completion.}
All cube systems use the same stopping check at 3~Hz. The agent must have bound the requested symbol to a goal, with no request pending and no inspection underway. Measured orientation error must be at most $\epsilon=0.16$ radians. This check can stop the system without an explicit stop action from the agent. On other objects, REVEAL uses calibrated completion correction, while the baselines use a measured-pose check. Their stopping rules are detailed in the supplement, and all are scored with the same true-pose tolerance. At stop time $\tau$, success is
\begin{equation}
 S=\mathbf{1}\!\left[\hat b_\tau=b^*(\ell),\ d(q_\tau,Q(b^*))\leq\epsilon,\ \neg D_{0:\tau}\right],
\label{eq:success}
\end{equation}
where $\hat b_\tau$ is the declared target binding, $q_\tau$ is true orientation, $D_{0:\tau}$ indicates an earlier terminal drop, and $d(q,Q)=2\arccos(|\langle q,Q\rangle|)$ for normalized quaternions. The target must be correct at the stop time, not merely reached earlier. We do not require an additional hold after stopping. A wrong binding or an incorrect final pose counts as failure even if the object remains in the hand. A drop terminates the trial when the object is at least 0.24 metres from its nominal grasp position.

\paragraph{Paired trials and service failures.}
Unless stated otherwise, each condition uses 40 trials: 20 per seed under two seeds. Compared methods share the same initial physical states and scored targets. Control and feedback experiments also share instructions. In the mapping experiment, the requested letter changes with the layout so that the physical target stays the same. A service error or a model call taking more than 120 seconds excludes that trial from the service-eligible comparison. We remove the same trial from every compared method, retain all raw results, and do not replace excluded trials. Each reported rate uses the denominator for its own comparison. Service failures are reported separately from control failures.
